\documentclass[a4paper,10pt]{article}
\usepackage[utf8]{inputenc}
\usepackage[T1]{fontenc}
\usepackage{lmodern}
\usepackage[ngerman]{babel}
\usepackage[landscape,margin=1.5cm]{geometry}
\usepackage{longtable,array,booktabs}
\usepackage{enumitem}
\usepackage{microtype}
\usepackage{fancyhdr}
\pagestyle{fancy}
\fancyhf{}
\fancyfoot[L]{\small\texttt{C3\_SPM\_Blatt2.pdf}}
\fancyfoot[R]{\small Seite \thepage}
\renewcommand{\headrulewidth}{0pt}
\renewcommand{\footrulewidth}{0.4pt}
\setlength{\footskip}{22pt}
\setlength{\parindent}{0pt}
\setlength{\parskip}{5pt}
\setlength{\emergencystretch}{3em}
\setlist{nosep}
\setlength{\tabcolsep}{4pt}
\renewcommand{\arraystretch}{1.15}
\setlength{\LTpre}{8pt}
\setlength{\LTpost}{12pt}
\newcommand{\sectionspace}[1]{\par\ifdim\dimexpr\pagegoal-\pagetotal\relax<#1\newpage\fi\addvspace{14pt}}
\title{MoinMetrics\\C3 SPM Blatt 2\\Aufgabe 1: Anforderungsanalyse}
\author{Gruppe C3 / MoinMetrics\\{}Alexander Borrmann\\{}Lukas Pancratius\\{}Ino Wölke}
\date{Hochschule Emden-Leer\\Software-Projektmanagement}
\begin{document}
    \maketitle
    \thispagestyle{fancy}
    \clearpage
    \subsection*{Bedingungen}
    \begin{itemize}
        \item Der Quellcode des Produkts wird mit ausgeliefert.
        \item Die Anwendung wird in Java programmiert.
        \item Die Analyse der Daten soll durch die Data-Mining Software Weka erfolgen.
        \item Es wird eine war-Datei ausgeliefert.
        \item Die Software wird in javadoc kommentiert.
        \item Alle sonstigen Dokumente liegen im PDF-Format vor.
        \item Das Front-End der Analysesoftware soll eine Webseite sein, die auf dem firmeninternen TomCat Webserver laufen soll.
        \item Die Webseite wird hauptsächlich durch die jeweilige Marktleitung (keine IT-Fachleute) benutzt werden.
        \item Die Anwendung muss durch ein Passwort geschützt werden, jedoch ist keine Benutzerverwaltung nötig.
    \end{itemize}

    \subsection*{Muss-Anforderungen an die Webseite}
    \begin{itemize}
        \item csv-Datei hochladen können, Fehler dürfen nicht zum Systemabsturz führen
        \item Graphische Aufbereitung der Daten
        \item Liste für allgemeine Marketingmaßnahmen, diese muss erweiterbar sein
    \end{itemize}

    \subsection*{Pflichtanforderungen}
    \begin{enumerate}
        \item \#1 10.000 Datensätze in 1-2 Sekunden, maximal 5 Sekunden
        \item \#12 Darstellung von Waren, die zusammen gekauft wurden (Apriori)
        \item \#5 Fehler in der csv-Datei sollten zu einer Fehlermeldung führen
        \item \#6 Umsatzstärkster Einkaufstag, Umsatzstärkste Einkaufsuhrzeit
        \item \#7 Stärkster Einkaufstag, stärkste Einkaufsuhrzeit, nach Kundenhäufigkeit
        \item \#8 Umsatzstärkste Einkaufsuhrzeit an den einzelnen Tagen
    \end{enumerate}

    \subsection*{optional}
    \begin{enumerate}
        \item \#9 Top und Flop - Artikel, nach Umsatz
        \item \#10 Top und Flop - Artikel, nach Deckungsbeitrag
    \end{enumerate}

    \clearpage
    \section*{Aufgabe 1: Anforderungsanalyse}

    Der folgende Katalog strukturiert die oben ausgewählten Anforderungen. Die Bezeichner F, Q und R stehen für Funktionsanforderungen, Qualitätsmerkmale und Randbedingungen. Die ursprünglichen Nummern aus der PDF bleiben zur Rückverfolgbarkeit angegeben.

\subsection*{a) Funktionsanforderungen}


    \normalsize
\begin{longtable}{>{\raggedright\arraybackslash}p{\dimexpr 0.085\linewidth-2\tabcolsep\relax}>{\raggedright\arraybackslash}p{\dimexpr 0.15\linewidth-2\tabcolsep\relax}>{\raggedright\arraybackslash}p{\dimexpr 0.365\linewidth-2\tabcolsep\relax}>{\raggedright\arraybackslash}p{\dimexpr 0.1\linewidth-2\tabcolsep\relax}>{\raggedright\arraybackslash}p{\dimexpr 0.3\linewidth-2\tabcolsep\relax}}
        \toprule
        \textbf{Bezeichner} & \textbf{Funktion} & \textbf{Beschreibung} & \textbf{Priorität / Bezug} & \textbf{Offene Abstimmungspunkte} \\
        \midrule
        \endhead
\bottomrule\endfoot
        F01 & CSV-Import & Die Anwendung muss eine CSV-Datei im vereinbarten Format hochladen und die enthaltenen Einkaufsdaten zur Analyse einlesen können. & Muss & - \\
        \midrule
        F02 & Fehlermeldung bei CSV-Fehlern & Die Anwendung muss bei einer fehlerhaften CSV-Datei eine Fehlermeldung anzeigen. & Pflicht, \#5 & \begin{minipage}[t]{\linewidth}\begin{itemize}[leftmargin=*,label=-]
\item Fehlerfälle
\item Umgang mit teilweise gültigen Dateien
            \end{itemize}\end{minipage} \\
        \midrule
        F03 & Passwortschutz & Die Anwendung muss vor dem Zugriff auf Daten und Auswertungen die Eingabe eines gültigen Passworts verlangen. & Muss & - \\
        \midrule
        F04 & Grafische Darstellung & Die Anwendung muss Analyseergebnisse grafisch darstellen. & Muss & \begin{minipage}[t]{\linewidth}\begin{itemize}[leftmargin=*,label=-]
\item Diagrammarten
\item Zuordnung zu Analysen
            \end{itemize}\end{minipage} \\
        \midrule
        F05 & Marketingmaßnahmen anzeigen & Die Anwendung muss eine Liste allgemeiner Marketingmaßnahmen anzeigen. \textbf{???} & Muss &  \\
        \midrule
        F06 & Marketingmaßnahmen erweitern & Die Liste allgemeiner Marketingmaßnahmen muss um zusätzliche Einträge erweitert werden können. & Muss & \begin{minipage}[t]{\linewidth}\begin{itemize}[leftmargin=*,label=-]
\item Art der Erweiterung
\item Zuständige Personen
            \end{itemize}\end{minipage} \\
        \midrule
        F07 & Gemeinsam gekaufte Waren & Die Anwendung muss mithilfe des Apriori-Verfahrens Waren ermitteln und darstellen, die zusammen gekauft wurden. & Pflicht, \#12 & \begin{minipage}[t]{\linewidth}\begin{itemize}[leftmargin=*,label=-]
\item Apriori-Parameter
\item Umfang der Ergebnisdarstellung
            \end{itemize}\end{minipage} \\
        \midrule
        F08 & Umsatzstärkster Einkaufstag & Die Anwendung muss den Einkaufstag mit dem höchsten Umsatz ermitteln und anzeigen. & Pflicht, \#6 & \begin{minipage}[t]{\linewidth}\begin{itemize}[leftmargin=*,label=-]
\item Auswertungszeitraum
\item Tag: Datum oder Wochentag
\item Umgang mit Gleichständen
            \end{itemize}\end{minipage} \\
        \midrule
        F09 & Umsatzstärkste Einkaufsuhrzeit & Die Anwendung muss die Einkaufsuhrzeit mit dem höchsten Umsatz ermitteln und anzeigen. & Pflicht, \#6 & \begin{minipage}[t]{\linewidth}\begin{itemize}[leftmargin=*,label=-]
\item Auswertungszeitraum
\item Zeitintervalle
\item Umgang mit Gleichständen
            \end{itemize}\end{minipage} \\
        \midrule
        F10 & Stärkster Einkaufstag nach Kundenhäufigkeit & Die Anwendung muss den Einkaufstag mit der höchsten Kundenhäufigkeit ermitteln und anzeigen. & Pflicht, \#7 & \begin{minipage}[t]{\linewidth}\begin{itemize}[leftmargin=*,label=-]
\item Auswertungszeitraum
\item Tag: Datum oder Wochentag
\item Umgang mit Gleichständen
\item Kundenhäufigkeit: unterschiedliche Kund*innen oder Einkaufsvorgänge
            \end{itemize}\end{minipage} \\
        \midrule
        F11 & Stärkste Einkaufsuhrzeit nach Kundenhäufigkeit & Die Anwendung muss die Einkaufsuhrzeit mit der höchsten Kundenhäufigkeit ermitteln und anzeigen. & Pflicht, \#7 & \begin{minipage}[t]{\linewidth}\begin{itemize}[leftmargin=*,label=-]
\item Auswertungszeitraum
\item Zeitintervalle
\item Umgang mit Gleichständen
\item Kundenhäufigkeit: unterschiedliche Kund*innen oder Einkaufsvorgänge
            \end{itemize}\end{minipage} \\
        \midrule
        F12 & Umsatzstärkste Einkaufsuhrzeit je Tag & Die Anwendung muss für jeden betrachteten Tag die Einkaufsuhrzeit mit dem höchsten Umsatz ermitteln und anzeigen. & Pflicht, \#8 & \begin{minipage}[t]{\linewidth}\begin{itemize}[leftmargin=*,label=-]
\item Auswertungszeitraum
\item Zeitintervalle
\item Tag: Datum oder Wochentag
\item Umgang mit Gleichständen
            \end{itemize}\end{minipage} \\*
        \midrule\noalign{\penalty10000}
        F13 & Top- und Flop-Artikel nach Umsatz & Die Anwendung soll die Artikel mit dem höchsten und niedrigsten Umsatz ermitteln und anzeigen. & Optional, \#9 & \begin{minipage}[t]{\linewidth}\begin{itemize}[leftmargin=*,label=-]
\item Auswertungszeitraum
\item Umgang mit Gleichständen
\item Anzahl der Top-/Flop-Ergebnisse
\item Artikel oder Warengruppen
            \end{itemize}\end{minipage} \\*
        \midrule\noalign{\penalty10000}
        F14 & Top- und Flop-Artikel nach Deckungsbeitrag & Die Anwendung soll die Artikel mit dem höchsten und niedrigsten Deckungsbeitrag ermitteln und anzeigen. & Optional, \#10 & \begin{minipage}[t]{\linewidth}\begin{itemize}[leftmargin=*,label=-]
\item Auswertungszeitraum
\item Umgang mit Gleichständen
\item Anzahl der Top-/Flop-Ergebnisse
\item Artikel oder Warengruppen
            \end{itemize}\end{minipage} \\
        \midrule
    \end{longtable}
\normalsize

\clearpage
\subsection*{b) Qualitätsmerkmale und Quantifizierung}


    \normalsize
\begin{longtable}{>{\raggedright\arraybackslash}p{\dimexpr 0.085\linewidth-2\tabcolsep\relax}>{\raggedright\arraybackslash}p{\dimexpr 0.14\linewidth-2\tabcolsep\relax}>{\raggedright\arraybackslash}p{\dimexpr 0.31\linewidth-2\tabcolsep\relax}>{\raggedright\arraybackslash}p{\dimexpr 0.205\linewidth-2\tabcolsep\relax}>{\raggedright\arraybackslash}p{\dimexpr 0.26\linewidth-2\tabcolsep\relax}}
        \toprule
        \textbf{Bezeichner} & \textbf{Qualitätsmerkmal} & \textbf{Messbares Kriterium} & \textbf{Nachweis / Status} & \textbf{Offene Abstimmungspunkte} \\
        \midrule
        \endhead
\bottomrule\endfoot
        Q01 & Analysezeit & Die Analyse von 10.000 Datensätzen muss spätestens nach 5 Sekunden abgeschlossen sein; 1-2 Sekunden sind der Zielwert. & Bezug: \#1. Vorschlag: Messung vom Start der Analyse bei bereits eingelesenen Daten bis zur vollständig sichtbaren Ergebnisausgabe; jede vereinbarte Analyse muss in 10 Testläufen die 5-Sekunden-Grenze einhalten. & \begin{minipage}[t]{\linewidth}\begin{itemize}[leftmargin=*,label=-]
\item Referenzhardware und Testdatei
\item Analyseparameter und einbezogene Analysen
\item Messbeginn und Messende
\item Messvorschlag und Anzahl der Testläufe bestätigen
            \end{itemize}\end{minipage} \\
        \midrule
        Q02 & Robustheit beim CSV-Import & Bei 100 \% der vereinbarten fehlerhaften CSV-Testdateien darf die Anwendung nicht abstürzen. Sie muss eine Fehlermeldung anzeigen und anschließend einen erneuten Upload ermöglichen. & Vorschlag zur Quantifizierung. Testfälle mindestens: falsches Trennzeichen, fehlende Pflichtspalten, ungültige Zahlen-/Datumswerte und leere Datei. & \begin{minipage}[t]{\linewidth}\begin{itemize}[leftmargin=*,label=-]
\item CSV-Format
\item Fehlerfallkatalog
\item Quantifizierung bestätigen
            \end{itemize}\end{minipage} \\
        \midrule
        Q03 & Bedienbarkeit & \textbf{???} & Vorschlag zur Quantifizierung für die Zielgruppe Marktleitung. & \begin{minipage}[t]{\linewidth}\begin{itemize}[leftmargin=*,label=-]
\item Quantifizierung und Testablauf
            \end{itemize}\end{minipage} \\*
        \midrule\noalign{\penalty10000}
        Q04 & Zugriffsschutz & Bei 100 \% der vereinbarten Zugriffstests ohne gültige Anmeldung dürfen weder Einkaufsdaten noch Analyseergebnisse ausgeliefert werden. Mit gültigem Passwort muss der Zugriff möglich sein. & Vorschlag zur Quantifizierung von F03. Tests umfassen den Aufruf der Oberfläche und direkte Aufrufe der datenliefernden Endpunkte. & \begin{minipage}[t]{\linewidth}\begin{itemize}[leftmargin=*,label=-]
\item Testumfang
\item Sitzungsverhalten
\item Quantifizierung bestätigen
            \end{itemize}\end{minipage} \\
        \midrule
    \end{longtable}
\normalsize

\clearpage
\subsection*{c) Randbedingungen}


    \normalsize
\begin{longtable}{>{\raggedright\arraybackslash}p{\dimexpr 0.085\linewidth-2\tabcolsep\relax}>{\raggedright\arraybackslash}p{\dimexpr 0.17\linewidth-2\tabcolsep\relax}>{\raggedright\arraybackslash}p{\dimexpr 0.47\linewidth-2\tabcolsep\relax}>{\raggedright\arraybackslash}p{\dimexpr 0.275\linewidth-2\tabcolsep\relax}}
        \toprule
        \textbf{Bezeichner} & \textbf{Merkmal} & \textbf{Beschreibung} & \textbf{Offene Abstimmungspunkte} \\
        \midrule
        \endhead
\bottomrule\endfoot
        R01 & Programmiersprache & Die Anwendung muss in Java implementiert werden. & \begin{minipage}[t]{\linewidth}\begin{itemize}[leftmargin=*,label=-]
\item Java-Version
            \end{itemize}\end{minipage} \\
        \midrule
        R02 & Analysewerkzeug & Die Datenanalyse muss mit der Data-Mining-Software Weka erfolgen. & \begin{minipage}[t]{\linewidth}\begin{itemize}[leftmargin=*,label=-]
\item Weka-Version
            \end{itemize}\end{minipage} \\
        \midrule
        R03 & Weboberfläche und Betrieb & Das Front-End muss als Webseite auf dem firmeninternen Tomcat-Webserver betrieben werden können. & \begin{minipage}[t]{\linewidth}\begin{itemize}[leftmargin=*,label=-]
\item Tomcat-Version
\item Zielumgebung
            \end{itemize}\end{minipage} \\
        \midrule
        R04 & Auslieferungsformat & Die Anwendung muss als WAR-Datei ausgeliefert werden. & - \\
        \midrule
        R05 & Quellcodeauslieferung & Der Quellcode des Produkts muss mit ausgeliefert werden. & - \\
        \midrule
        R06 & Quellcodedokumentation & Die Software muss mit Javadoc kommentiert werden. Vorschlag: alle öffentlichen Klassen und Methoden. & \begin{minipage}[t]{\linewidth}\begin{itemize}[leftmargin=*,label=-]
\item Umfang der Javadoc-Kommentierung
            \end{itemize}\end{minipage} \\
        \midrule
        R07 & Dokumentformate & Alle sonstigen Dokumente müssen im PDF-Format vorliegen. & - \\*
        \midrule\noalign{\penalty10000}
        R08 & Eingabedaten & Die Einkaufsdaten müssen aus CSV-Dateien im festgelegten Format übernommen werden. & - \\*
        \midrule\noalign{\penalty10000}
        R09 & Umfang der Benutzerverwaltung & Eine Benutzerverwaltung ist nicht erforderlich und gehört nicht zum vereinbarten Funktionsumfang. Der Passwortschutz nach F03 bleibt erforderlich. & - \\
        \midrule
    \end{longtable}
\normalsize

\sectionspace{180pt}
\subsection*{d) Eigenschaften gut formulierter Anforderungen}

Eine Anforderungsspezifikation soll folgende Eigenschaften erfüllen:

\begin{itemize}
\item \textbf{Eindeutig:} Die sprachliche Formulierung lässt nur eine Interpretation zu. Unklare Begriffe wie „gut“ oder „schnell“ werden konkretisiert.
\item \textbf{Vollständig:} Die Spezifikation enthält alle erforderlichen Anforderungen und alle Bestandteile, die der zugrunde gelegte Standard verlangt.
\item \textbf{Überprüfbar:} Für jede Forderung lässt sich durch einen Test oder eine Prüfung feststellen, ob sie erfüllt ist. Qualitätsmerkmale erhalten dazu messbare Kriterien und definierte Testbedingungen.
\item \textbf{Konsistent:} Begriffe werden einheitlich verwendet; Synonyme für denselben Sachverhalt werden vermieden. Die Anforderungen stimmen mit den geltenden Standards überein und enthalten keine logischen, zeitlichen oder ursächlichen Widersprüche.
\item \textbf{Änderbar:} Eine klare Gliederung, ein Inhaltsverzeichnis, ein Index und Querverweise erleichtern Änderungen. Redundante Beschreibungen werden vermieden, damit eine Änderung nicht an mehreren Stellen nachgeführt werden muss.
\item \textbf{Verfolgbar:} Entscheidungen und ihre Begründungen sind nachvollziehbar dokumentiert. Eindeutige Bezeichner und Verweise ermöglichen die Zuordnung der Anforderungen zu ihren Quellen, ihrer Umsetzung und ihren Tests.
\end{itemize}
\end{document}
